\section{Propuesta de Adecuación y Equipamiento}
\label{sec:adecuacion}

\subsection{Obras y Reformas Propuestas}
Para optimizar el flujo de usuarios y mejorar el confort térmico y acústico de las cafeterías, se proponen las siguientes actuaciones de adecuación física:
\begin{enumerate}
    \item Renovación de revestimientos y pavimentos con materiales antideslizantes y de fácil limpieza.
    \item Mejora del sistema de climatización y ventilación.
    \item Rediseño de la línea de self-service para agilizar los tiempos de espera en las horas punta.
\end{enumerate}

\subsection{Equipamiento y Mobiliario}
Se renovará la maquinaria de cocina y barra, priorizando equipos de alta eficiencia energética para reducir la huella de carbono del centro hospitalario. En la Tabla~\ref{tab:equipos} se resume el equipamiento principal a incorporar.

\begin{table}[htbp]
    \centering
    \caption{Relación de equipamiento principal propuesto.}
    \label{tab:equipos}
    \begin{tabular}{llcc}
        \toprule
        \textbf{Descripción} & \textbf{Ubicación} & \textbf{Unidades} & \textbf{Clase Energética} \\
        \midrule
        Lavavajillas de cúpula industrial & Cocina Principal & 2 & A+++ \\
        Vitrina expositora refrigerada & Barra Cafetería & 3 & A+ \\
        Cafetera exprés profesional & Barra Cafetería & 2 & A \\
        Hornos de convección mixtos & Cocina Principal & 1 & A++ \\
        \bottomrule
    \end{tabular}
\end{table}

\subsection{Puntos de Venta Automática (Vending)}
Los puntos de vending serán renovados por máquinas inteligentes (smart vending) con lector de tarjetas bancarias, dispositivos móviles y tarjetas de empleado del HUVR. Al menos el $30\%$ de los productos ofrecidos corresponderán a la gama de alimentación saludable (fruta fresca, frutos secos naturales, lácteos desnatados, bebidas sin azúcares añadidos), cumpliendo las directrices de la Consejería de Salud y Consumo \cite{salud_andalucia}.
